% ============================================================
% 1  Mermin-Wagner and the ferromagnetic loophole
% ============================================================
\begin{frame}[t]
\ft{No continuous order in one dimension}\par
\vspace{0.3cm}
\begin{columns}[c]
\begin{column}{0.34\textwidth}
\centering
% order parameter in the complex plane: amplitude fixed on the ring, phase free
\begin{tikzpicture}[scale=0.8,font=\scriptsize]
 \draw[-{Stealth[length=3pt]},mygray] (-1.25,0)--(1.35,0) node[right]{Re\,$\psi$};
 \draw[-{Stealth[length=3pt]},mygray] (0,-1.2)--(0,1.3) node[above]{Im\,$\psi$};
 \draw[pumpkin,line width=1.6pt] (0,0) circle (0.95);
 \draw[-{Stealth[length=5pt]},very thick,cLL] (0,0)--(40:0.95);
 \draw[cLL] (0.35,0) arc (0:40:0.35);
 \node[cLL] at (20:0.55) {$\theta$};
 \node[pumpkin!80!black,anchor=west] at (-62:1.0) {$|\psi|=\sqrt{\rho_0}$};
\end{tikzpicture}\par
\end{column}
\begin{column}{0.64\textwidth}
{\small U(1) order parameter\ \ $\psi=\colorbox{pumpkin!50}{$\mystrut\sqrt{\rho_0}$}\,e^{i{\color{cLL}\theta}}$\par}
\vspace{0.1cm}
{\footnotesize ordered: \colorbox{pumpkin!50}{\mystrut amplitude} fixed by the condensate,\\ only the \colorbox{cLL!20}{\mystrut phase} can vary\par}
\vspace{0.2cm}
\onslide<2->{%
\begin{tikzpicture}[x=0.42cm]
 \foreach \i in {0,...,13} {
   \pgfmathsetmacro{\a}{40+55*sin(24*\i)}
   \fill[mygray] (\i,0) circle (0.04);
   \draw[-{Stealth[length=4.5pt,width=4pt]},thick,cLL] (\i,0) -- ++(\a:0.5);
 }
 \node[font=\scriptsize,anchor=west] at (13.5,0.15) {$\theta(x,\tau)$};
\end{tikzpicture}\par}
\end{column}
\end{columns}
\vspace{0.1cm}
\onslide<2->{%
\centering
{\small$\displaystyle S=\frac{\rho_s}{2}\!\int\!\dd\tau\,\dd x\,\Big[(\partial_\tau\theta)^2+(\partial_x\theta)^2\Big]
\quad\Rightarrow\quad \langle\theta\theta\rangle_{\omega,k}=\frac{1}{\rho_s(\omega^2+k^2)}$\par}
\vspace{0.05cm}
{\scriptsize lowest order in gradients allowed by the symmetry \colorbox{citron!40}{\mystrut $\theta\to\theta+\text{const}$}\par}}
\vspace{0.15cm}
\begin{columns}[T]
\begin{column}{0.48\textwidth}
\centering
\onslide<3->{\colorbox{mygraylight!60}{\parbox{0.92\linewidth}{\centering\footnotesize
 \textbf{2+1D}\\[1pt]
 $\langle(\theta_x-\theta_0)^2\rangle\propto\int\!\dd^2k\,\dd\omega\,\frac{1-\cos k\cdot x}{\omega^2+k^2}$\\[3pt]
 \onslide<4->{$=\int\!\dd^2k\,\frac{1-\cos k\cdot x}{2|k|}$}\\[3pt]
 \onslide<5->{$\sim\int\!\frac{k\,\dd k}{k}$\ \ \textbf{finite}}\\[4pt]
 \onslide<5->{$\langle e^{i\theta_x}e^{-i\theta_0}\rangle\to$ const:\ order survives}}}}
\end{column}
\begin{column}{0.48\textwidth}
\centering
\onslide<6->{\colorbox{cLL!20}{\parbox{0.92\linewidth}{\centering\footnotesize
 \textbf{1+1D}\\[1pt]
 $\langle(\theta_x-\theta_0)^2\rangle\propto\int\!\dd k\,\dd\omega\,\frac{1-\cos kx}{\omega^2+k^2}$\\[3pt]
 \onslide<7->{$=\int\!\dd k\,\frac{1-\cos kx}{2|k|}$}\\[3pt]
 \onslide<8->{$\sim\int_{1/|x|}^{\Lambda}\!\frac{\dd k}{k}=\ln\Lambda|x|$\ \ \textbf{diverges}}\\[4pt]
 \onslide<8->{$\langle e^{i\theta_x}e^{-i\theta_0}\rangle\propto|x|^{-\eta}\to0$:\ no order}}}}
\end{column}
\end{columns}
\vspace{0.05cm}
{\tiny\color{citegray} N.~D.~Mermin, H.~Wagner, PRL 17, 1133 (1966);\; P.~C.~Hohenberg, PR 158, 383 (1967);\; S.~Coleman, CMP 31, 259 (1973)\par}
\onslide<8->{\tk{In 1+1D, phase fluctuations destroy order (Mermin--Wagner)}}
\end{frame}
% Notes: Step 1: A U(1) order parameter is a complex number: an amplitude and a phase (superfluid, XY magnet). Suppose we are ordered: the amplitude is pinned by the condensate density rho_0 (the orange ring), but the phase theta can point anywhere on the ring: all directions cost the same.
% Notes: Step 2: So the phase can vary in space and time (the arrows). Write the cheapest action that respects theta -> theta + const: no theta^2 term allowed, only gradients; lowest order is one time and one space derivative squared. Its Green's function is 1/(omega^2 + k^2).
% Notes: Steps 3-5 (2+1D first, the familiar case): how much does the phase wander between two distant points? Do the frequency integral: 1/|k|. In two dimensions the measure d^2k = k dk eats the 1/k: the integral is finite. The phase stays put on average, <e^{i theta} e^{-i theta}> -> const: a superfluid / XY magnet orders.
% Notes: Steps 6-8 (1+1D): same steps, but now the measure is only dk. The 1-cos(kx) cuts the integral off at k ~ 1/|x|: we get ln(Lambda |x|), growing without bound. Then <e^{i theta} e^{-i theta}> = exp(-1/2 <dtheta^2>) ~ |x|^{-eta}: power law, order is destroyed. This is Mermin-Wagner (Coleman for 1+1D quantum).


\begin{frame}[t]
\ft{The ferromagnetic loophole}
\begin{columns}[T]
\begin{column}{0.48\textwidth}
\centering
\vspace{0.15cm}
$\displaystyle \mathcal L = \colorbox{orchid!30}{$\mystrut i\rho\,\partial_\tau\theta$} + \frac{\rho^2}{2\chi} + \frac{\rho_s}{2}(\partial_x\theta)^2$\\[4pt]
{\scriptsize\colorbox{orchid!30}{\mystrut Berry phase}: density $\rho=S^z$ is conjugate to $\theta$;\\ first order in time, not a classical weight\par}
\end{column}
\begin{column}{0.48\textwidth}
\begin{itemize}
 \item Ferromagnets, bosons: \colorbox{citron!40}{\mystrut Berry phase}
 \onslide<2->{\item At the critical point $1/\chi\to0$:\\ density cost $\rho^2\to(\partial_x\rho)^2$}
\end{itemize}
\end{column}
\end{columns}
\vspace{0.45cm}
\begin{columns}[T]
\begin{column}{0.48\textwidth}
\centering
\colorbox{cLL!20}{\parbox{0.92\linewidth}{\centering\footnotesize
 \textbf{\color{cLL}Luttinger liquid}\ \ $\omega\propto k$,\ $z=1$\\[3pt]
 $\displaystyle\langle\delta\theta^2\rangle\propto\int\!\dd k\,\dd\omega\,\frac{1}{\omega^2+k^2}\sim\int\!\frac{\dd k}{k}=\infty$}}
\end{column}
\begin{column}{0.48\textwidth}
\centering
\onslide<2->{\colorbox{cQCP!15}{\parbox{0.92\linewidth}{\centering\footnotesize
 \textbf{\color{cQCP}Critical point}\ \ $\omega\propto k^2$,\ $z=2$\\[3pt]
 $\displaystyle\langle\delta\theta^2\rangle\propto\int\!\dd k\,\dd\omega\,\frac{k^2}{\omega^2+k^4}\sim\int\!\dd k\ \ \text{finite}$}}}
\end{column}
\end{columns}
\vspace{0.45cm}
\onslide<3->{%
\centering
{\footnotesize\colorbox{pumpkin!50}{\mystrut For the expert}\ \ $K(\partial_\tau\phi)^2$ not renormalized $\Rightarrow[\partial_x\phi]=\tfrac{3-z}{2}$;\ \ unbroken U(1): $[\rho]=1$\\[3pt]
$\Rightarrow$\ \colorbox{cQCP!15}{\mystrut $z\neq1$ forces U(1) breaking}\par}
\vspace{0.1cm}
{\tiny\color{citegray} M.~P.~A.~Fisher, G.~Grinstein, S.~M.~Girvin, PRL 64, 587 (1990);\; X.-G.~Wen, PRB 46, 2655 (1992);\; S.~Sachdev, Z.~Phys.~B 94, 469 (1994)\par}}
\tk{Soft density $\Rightarrow$ stiff phase}
\end{frame}
% Notes: Ferromagnets have a Berry phase: the magnetization density rho=S^z is the momentum conjugate to the angle theta. This term is first order in imaginary time, so the Euclidean weight is complex: no classical stat-mech model to map to.
% Notes: Berry phase alone is not enough (left box): with a finite susceptibility we are back to omega ~ k and log fluctuations.
% Notes: Step 2: at the ferromagnetic critical point the uniform susceptibility diverges; the density costs only gradients. Integrating out rho now gives theta a k^2/(omega^2+k^4) propagator: the infrared integral converges. Phase fluctuations are finite: true order is allowed.
% Notes: Step 3 (experts): nonperturbative version. Interactions only involve gradients of phi, so the (d_tau phi)^2 coefficient is not renormalized; this fixes [d_x phi]=(3-z)/2. A conserved, unbroken U(1) density must have dimension 1, which needs z=1. Any z != 1 forces symmetry breaking.
